\section{Question: what exactly is FFGFT's time-winding?}
In FFGFT, time is the mass circle, woven into the remaining circles via $\tilde T\cdot m=1$ and
not a separable factor (Dok.~285). A precise question follows: what exactly \emph{is} FFGFT's
time-winding?

This document settles it completely FFGFT-internally --- it is \emph{derivable} from the fractal
correction. The result in one sentence: the K$_{\text{frak}}$ recursion forces the time-winding
into a \emph{logarithmic spiral with self-similarity ratio $e$}, that is, into continuous scale
invariance.

\section{Three building blocks from the corpus}
Three quantities are fixed elsewhere and are only brought together here. First, the fractal
dimension of space, $D_f=3-\xi_0$ with $\xi_0=4/30000$ (Dok.~205, Dok.~101). Second, the fractal
correction in the gate: the X rotation is not by $\pi$ but by $\alpha=\pi\,K_{\text{frak}}$ with
$K_{\text{frak}}=1-100\,\xi_0$ (Dok.~034). Third, the scale run as the recursion
$\xi_{n+1}=\xi_n\,(1-100\,\xi_0)$, more precisely $\xi_{n+1}=\xi_n(1-100\,\xi_n)$ (Dok.~146).

From $\xi_0=1/7500$ follow three exact facts that carry the rest:
\begin{equation}
  100\,\xi_0=\frac{1}{75},\qquad K_{\text{frak}}=\frac{74}{75},\qquad
  \frac{1}{100\,\xi_0}=75 .
\end{equation}

\section{Non-closure: from the 75-closure to the spiral}
A winding that should advance by $2\pi$ per turn advances gate-wise only by
$\pi\,K_{\text{frak}}$ and \emph{misses} $\pi$ by $100\pi\,\xi_0$ per half-turn. Iterated, the
winding does not close but precesses --- it becomes a spiral. The closure defect per turn is
$d=100\,\xi_0$ in turn units.

For frozen $\xi_0$, $d=1/75$ is constant, so the rotation number $74/75$ is rational: the winding
closes after \emph{exactly $75$ turns} and carries a $75$-fold structure --- the number
$75=1/(100\,\xi_0)$ drops out exactly. The limiting case of constant pitch, $d=0$, closes
already after one turn. Both closed cases are the
\emph{rolled-up}, compact reading of the time-winding: a single scale, the winding lies closed
and is static. FFGFT already knows a non-closure
algebraically, in the non-multiplicative mode composition (Dok.~193); here its
geometric-temporal counterpart is added.

\section{The scale run makes the spiral logarithmic}
When $\xi$ runs over the recursion, the rational $75$-closure turns into a non-closing,
self-similar spiral. This transition is the \emph{unrolling} of the winding: the spiral is the
decompactified counterpart of the closed $75$-winding, and \enquote{unrolling} here means
precisely letting the scale run (the recursion), not merely splitting off a torus factor. The
derivation is closed. With $u_n=1/\xi_n$,
\begin{equation}
  u_{n+1}=\frac{1}{\xi_n(1-100\,\xi_n)}=u_n\,(1+100\,\xi_n+\mathcal{O}(\xi_n^2))
         =u_n+100+\mathcal{O}(\xi_n),
\end{equation}
so $u_n\approx 100\,n+7500$ and hence an \emph{algebraic} $1/n$ decay
\begin{equation}
  \xi_n\;\approx\;\frac{1}{100\,(n+75)} .
\end{equation}
The cumulative precession defect thereby becomes logarithmic,
\begin{equation}
  D(k)=\sum_{n\le k}100\,\xi_n\;\approx\;\sum_{n\le k}\frac{1}{n+75}
       \;\approx\;\ln\!\Big(1+\frac{k}{75}\Big).
\end{equation}
With radius $\rho=k+75$ --- the pole of the spiral sits at the effective origin
$n=-75=-1/(100\,\xi_0)$, the same $75$ as before --- and precession angle $\beta=2\pi\,D(k)$, it
follows that $\rho/75=e^{\beta/2\pi}$, that is
\begin{equation}
  \rho\;=\;75\,e^{\beta/2\pi}\;=\;\text{logarithmic (equiangular) spiral.}
\end{equation}
Its self-similarity is thereby fixed: per precession round, $\Delta\beta=2\pi$, the radius scales
by $e^{1}=e$. The \emph{ratio $e$ is not assumed but follows} from the algebraic $1/n$ decay of
the recursion: it is the base of the natural logarithm produced by the harmonic sum, and no
additional parameter enters. The radius offset $75$ is likewise not freely chosen --- it is
$1/(100\,\xi_0)$, the same constant that already carries the $75$-closure.

Both are confirmed numerically without leeway (\texttt{zeit\_logspirale\_probe.py}): the fit
$\ln\rho=a\,\beta+b$ gives $a=0{.}15922\approx 1/2\pi=0{.}15915$ and $e^{2\pi a}=2{.}7195\approx e$;
the locally measured self-similarity ratio converges scale by scale
$2{.}734\to2{.}726\to2{.}721\to2{.}720\to2{.}719\to2{.}719$ towards $e=2{.}71828$.

\section{Continuous, not discrete scale invariance}
The self-similarity is \emph{continuous} --- a log spiral, a power law --- not discretely
log-periodic. A test for discrete scale invariance (the residual of $D(k)$ against the smooth
$\ln(1+k/75)$) yields an amplitude near zero and no periodicity in $\ln k$. This agrees with the
corpus: Dok.~146 records that the scaling over the factor $1/\xi_0=7500$ becomes
quasi-continuous through about $13$ logarithmic sub-levels, and the fractal dimension stays
constant without steps. The time-winding realises the same continuous scale invariance
geometrically, as a log spiral.

\section{Space and time: the same $\xi_0$, scaled differently}
Space and time carry the same fundamental constant $\xi_0$ at two places: the spatial
box-dimension deficit is $\xi_0$ (in $D_f=3-\xi_0$), the winding defect per gate is
$100\,\xi_0=1/75$. The factor $100$ between them is the recursion coefficient; it links a
\emph{geometric} deficit to a \emph{dynamical} defect and privileges no axis. Which axis carries
the defect is a matter of projection --- it is conventionally noted on time, but is symmetric
under $\tilde T\cdot m=1$ (following sections), hence no amplification of time over space.

A boundary must be drawn explicitly here: the box dimension of the time-winding itself is
$\approx 1$ --- it is a curve --- and \emph{not} $3-\xi_0$. The link between the temporal spiral
and the spatial fractal dimension is the shared $\xi_0$, not an equality of dimensions. The time
spiral does not carry the space dimension.

\section{The choice of axis is a projection}
The logarithmic spiral is a property of the whole coupled winding on $T^4$, not of a single
coordinate. What is forced is the \emph{non-closure} with ratio $e$; onto which axis one places
it is a matter of representation. Because $\tilde T\cdot m=1$ couples time and mass inversely and
the torus weaves all circles together --- nothing factorises --- there is no distinguished time
dimension that rolls up on its own. A generic winding runs through several circles at once; the
pure \emph{time} winding is the special case of a projection. The same non-closure could equally
be placed on a space axis or on a mixture --- different basis, same invariant.

The assignment to time is thus an artificial shift onto one axis --- a convention, not a
distinction. The gate defect is conventionally described on the time-winding (Dok.~034); the dual
projection (next section) shows, however, that the same non-closure appears with the same factor
$100$ on the space side. No basis is thereby preferred; each is an equivalent representation of
the same invariant --- the same representational freedom as between the two equivalent
representations (Dok.~206) and in the translation invariance of $\theta=2/9$ (Dok.~293).

\section{Continuous time: the dual projection into space}
If, conversely, time is set \emph{continuous} --- smooth, without gate-wise rolling up --- then
the non-closure can no longer sit on time; it must be carried by another axis. The coupling
$\tilde T\cdot m=1$ does not allow time to be smoothed in isolation: with
$d\tilde T/\tilde T=-\,dm/m$ the mass runs along inversely, and with it the spatial-geometric side
($D_f=3-\xi_0$). The running winding then appears as a \emph{variation of space} rather than as
rolled-up time. It is the same invariant non-closure with ratio $e$, only rotated into the dual
projection: in the present reading time runs (rolls up) and space stays comparable; in the dual
reading time is continuous and space varies.

This mirroring has been worked out numerically (\texttt{dual\_projektion\_probe.py}). Because time
advances per gate shortened by the factor $K_{\text{frak}}=1-100\,\xi_n$, mass advances by the
inverse factor $1/K_{\text{frak}}$ under $\tilde T\cdot m=1$; the defect per step is $100\,\xi_n$
on both sides (the ratio $d_{\text{mass}}/d_{\text{time}}\to1$ to six digits). The factor $100$
therefore does not \emph{move}, it is \emph{mirrored}: the gate factors are inverse
($K_{\text{frak}}\leftrightarrow1/K_{\text{frak}}$), the amplification $100$ is the same. Both
projections yield the same logarithmic growth, and the spiral fit gives the same ratio $e$ for
both ($e^{2\pi a}=2{.}7195$ and $2{.}7192$); their cumulative defects differ only by a
\emph{bounded} constant ($\approx0{.}0134$, not growing with $\ln k$). The core content is thereby
representation-invariant and the non-closure preserved: time cannot be smoothed without mass and
space taking over the winding --- the mass-time duality principle in action. Its distribution over
time, mass and space is a matter of projection; the ratio $e$ is not.

\section{Placement relative to HLV}
The FFGFT-internal statement is thereby \emph{derived}: the K$_{\text{frak}}$ recursion forces the
time-winding to be a log spiral with ratio $e$, equivalent to a $1/n$ decay of its pitch. Only here
does HLV enter --- and here it must first be separated that \enquote{spiral time} in HLV denotes
not one object but two. First the \emph{geometric} time from the factorisation $T^7=T^4\times T^3$
(Dok.~285): there time is an $A_5$ singlet on the $S^1$ factor, decoupled and separable. Second the
\emph{temporal-dynamic} spiral time, the object HLV \emph{calls} by that name: a non-Markov
operator with a triadic time coordinate and a memory channel, without $A_5$ and without an $S^1$
factor. FFGFT, unlike both, does not factorise; its spiral is the projection of an inseparable
whole.

The limiting case of constant pitch ($d=0$) noted above corresponds to the fixed point $\xi\to0$ of
the recursion, in which the defect vanishes and the winding closes. That is the \emph{geometric}
decoupled singlet time, not the temporal-dynamic object; the running FFGFT winding approaches it
asymptotically. The decisive question to HLV is therefore no longer \enquote{similar?} but
\enquote{ratio $e$, yes or no?} --- and it must be put to the right one of the two objects. The
following section carries out this check. Ordered within the three methodological layers: the
derivation recursion~$\to$~log spiral is a core derivation, obtained from $\xi$; the identification
with one of the two HLV times is a bridge candidate; no physical claim over HLV is made.

\section{Comparison with HLV's spiral time: shared $Z_3$, no ratio $e$}
The identification of the FFGFT log spiral with HLV's \enquote{spiral time} can now be checked
against a concrete HLV source --- namely against the one of the two objects designated
\emph{temporal-dynamic} in the previous section. Formally it is a non-Markov operator with a triadic
time coordinate
\begin{equation}
  \Psi(t)=(t,\phi(t),\chi(t)),\qquad \Psi(t)=t+i\,\phi(t)+j\,\chi(t),
\end{equation}
where $t$ notes the laboratory ordering, $\phi$ the phase organisation and $\chi$ a memory; it
carries neither $A_5$ nor a singlet nor an $S^1$ factor and possesses a memory kernel $K_\chi$ with
Markov limit $g_\chi\to0$. Precisely this object HLV calls \emph{spiral time}; the decoupled $A_5$
singlet (previous sections) is the other. The two share only the name.

\section{The temporal-dynamic object does not carry the ratio $e$}
The test question posed above --- \enquote{ratio $e$, yes or no?} --- targets a geometrically
self-similar object. The temporal-dynamic object is not of that category. It is rather an object of
open quantum dynamics: its total dynamics is --- per Marcel's spiral-time paper --- a unitary
dilation on the product space $H_S\otimes H_M$ of system and memory, and it is tested there by
open-system criteria (process tensor, CP-divisibility, hard reset, finite-bath comparison), not by
a geometric self-similarity. It thereby belongs recognisably to the dynamical, not the geometric
category; the $e$-question is put to an object that, by its own construction, makes no such claim.
It carries no derived
self-similarity ratio, no $1/n$ decay of a pitch and no scale run; its non-closure sits in the
memory channel $\chi$, not in a running scale. Its only self-similar internal structure is the
golden-ratio / phason order of the associated cut-and-project (the control replaces the $\varphi$
structure by random irrational phases), that is $\varphi$ and not $e$. On the $e$-criterion the
identification is therefore not supported. What both frameworks structurally share is the role of
non-closure as the temporal core object and an explicit memory-off limit: FFGFT's $d=0$ closure
($\xi\to0$) corresponds to Marcel's Markov fallback $g_\chi\to0$. Same role, different mechanism
--- geometric accumulation ($D(k)=\sum 100\,\xi_n\to\ln$) against a stochastic memory kernel.

FFGFT's non-closure is available not only geometrically as a log spiral but equivalently in the
Hilbert space: as a memory kernel, the Fourier transform of the discrete spectral density
$\sum_k g_k^2\,\delta(\omega-\omega_k)$ of the $T^4$ connection-Laplacian (Dok.~283) --- the same
inseparable object in two representations ($H=L^2(T^4)\otimes\mathbb{C}^3$, a lossless bijection,
Dok.~230/282); this is representational robustness, not additional evidence. The comparison can
therefore be run in Marcel's own language, kernel against kernel: FFGFT's kernel is oscillatory and
discrete-spectral (revivals, BLP backflow $5.125$), Marcel's decaying and bounded --- the same
object class, a different form (Dok.~283). Precisely the $1/n$ scale run that carries $e$ is what
the decaying kernel lacks.

\section{Computational check of the $e$-criterion}
The $e$-criterion is necessary and sufficient via the pitch decay: a winding is equiangular (a log
spiral, ratio $e$ per $2\pi$) exactly when the defect per step decays like $1/n$ and the cumulative
defect grows logarithmically. A constant defect, by contrast, gives an Archimedean winding with
ratio $\to1$. Both constructions can be measured against this on their own definitions
(\texttt{spiralzeit\_verhaeltnis\_probe.py}).

FFGFT carries $e$: with $d_n=100\,\xi_n$ one has $d_n\,(n+75)\to1$ --- the $1/n$ decay --- and the
spiral fit gives $e^{2\pi a}=2{.}71947\approx e$. Marcel's spiral time does not: its phase
$\phi_i=q_\phi\,u_i$ is linear in the internal coordinate and its memory $\chi_i=F(\rho_i)$ is a
bounded window; the defect per step stays constant, $D(k)$ grows linearly rather than
logarithmically ($D(10^5)=1333$ against $\ln(1+k/75)=7{.}2$), and the measured ratio is $1{.}0011$.
On its own definitions it is thus not a log spiral at all, but a winding of constant pitch.

The geometry HLV actually carries points, moreover, to a constant other than $e$. Were one to build
an equiangular spiral from the $\tau=\varphi$ cut-and-project, its natural self-similarity factor
would be golden --- $\varphi=1{.}618$ per turn, or $\varphi^4=6{.}854$ per $2\pi$ --- not
$e=2{.}718$. Both frameworks thus know, in so far as they carry self-similarity at all, different
constants: $e$ on the FFGFT side (from the $1/n$ decay of the $\xi$ recursion), $\varphi$ on the
geometric HLV side (from the cut-and-project). As a bridge condition in the sense of P35 this reads
precisely: HLV's spiral time carries $e$ only if its phase law contains a $1/n$ scale run --- a
running internal scale ---; the given definitions do not contain it. That is a candidate, not a
derivation.

\section{Basis and dimensions: same numbers, different meaning}
Whether both frameworks carry the \emph{same} basis can be ordered dimension by dimension. The
numbers coincide, the meaning of the dimensions does not. The \emph{three}: FFGFT's three is the
internal $\mathbb{C}^3$, i.e. the $Z_3$ order on $T^4/Z_3$ (Dok.~282) --- a mode and generation
index over the torus. Marcel's three is the triadic \emph{time} $(t,\phi,\chi)$ --- a decomposition
of the one ordering axis into order, phase and memory. In his root-mass note this three coincides
with a $C_3/Z_3$ orbit ($2\pi k/3$, $k=0,1,2$); to that extent both frameworks share the $Z_3$
($C_3<A_5$, Dok.~285). What is shared is the $Z_3$, not a common three-dimensional coordinate
space: one indexes modes, the other time roles.

The \emph{six}: FFGFT's six is not native. It appears only via the HLV embedding
$T^7=T^4\times T^3$ (Dok.~285); FFGFT itself is four-dimensional ($T^4$). Marcel's six is the parent
lattice $\mathbb{Z}^6$ of the cut-and-project, split into three physical directions
($P_\parallel$) and three internal, phason ones ($P_\perp$). The same number, a different anatomy:
a $4+3$ torus embedding against a $3+3$ cut-and-project. Moreover, in Marcel's case the six is not
the basis of the \emph{continuous} spiral time --- that one is triadic (the three) --- but of the
\emph{discrete} $G$ lattice; the two levels are for him explicitly representations of the same
operator, not a replacement of it. This corresponds to the representational freedom on the FFGFT
side (two equivalent representations, Dok.~206; the rolled-up $75$-winding against the unrolled log
spiral).

The only concretely established overlap is the $\varphi$-icosahedral shared object ($C_3<A_5$),
computationally established but \emph{one-sided} --- the test direction is FFGFT$\to$HLV
(Dok.~294) --- and filtration-sensitive: it separates HLV sharply from crystalline controls but
only weakly from a generic $\sqrt2$ cut-and-project. Computationally this can be pinned to the axis
spectrum (\texttt{basis\_6d\_achsen\_probe.py}): the six projected axes of the $6$D cut-and-project
are equiangular only at $\tau=\varphi$ --- all fifteen pairwise angles equal to $63{.}435^\circ$ ---
while $\tau=\sqrt2$ gives a two-valued spectrum ($61{.}87^\circ/70{.}53^\circ$). The six therefore
do not coincide per se; what is equal is the golden-icosahedral projection alone. Basis equality is
thus a \emph{candidate}, not a derived identity. The numbers $3$ and $6$ coincide, the shared
invariant is the $Z_3$, and the dimensions themselves carry different meanings.

\section{Status}
Secured and derived: the time-winding as a logarithmic spiral with self-similarity ratio $e$,
continuous scale invariance, from the K$_{\text{frak}}$ recursion; the exact $75$-closure in the
frozen case; the shared constant $\xi_0$, appearing as the spatial box deficit $\xi_0$ and as the
winding defect $100\,\xi_0$ (factor $100$ = recursion coefficient, axis-symmetric). The closed
$75$-winding is here the rolled-up (compact, single-scale) reading, the log spiral the unrolled
(decompactified, across-scales) reading of the same time-winding.

The comparison with HLV's spiral time is negative on the $e$-criterion: the temporal-dynamic object
carries no ratio $e$, but golden-ratio / phason quasiperiodicity; the basis question reduces to the
shared $Z_3$ ($C_3<A_5$, one-sidedly verified, filtration-sensitive). The sharp residual question
shifts to the \emph{geometric} HLV object --- whether the factorised layer ($A_5$ singlet on $S^1$)
knows a scale-invariant object that would be testable against $e$ ---; in the available sources no
such object is exhibited. The category separation (two spiral times) and the dimension-by-dimension
comparison are secured findings; the identification remains a bridge candidate, currently blocked on
the $e$-criterion.

The contrast here is structural, not ontological. Both frameworks --- and likewise the two variants
within FFGFT, whether time or space is rolled up --- are mathematical models, not reality; no
factorisation and no choice of axis is thereby ontologically privileged. That one hopes such a
model is applicable to reality is the motivation, not a statement about what is.

Reproducible in \texttt{zeit\_windung\_probe.py}, \texttt{zeit\_logspirale\_probe.py},
\texttt{dual\_projektion\_probe.py}, \texttt{spiralzeit\_verhaeltnis\_probe.py} and
\texttt{basis\_6d\_achsen\_probe.py}.
